\subsection{Financial Control Extension to the Denial Attack Simulator}
\label{app:financial_control_simulator_update}

The D-ACT Denial Attack Simulator was extended with a victim financial profile and mitigation-specific financial impact model. The extension remains a purely visual and educational artefact. It does not generate network traffic, replay packet captures, invoke cloud services, trigger serverless functions, or consume AI APIs. Its purpose is to demonstrate how denial attacks may translate into financial exposure for a victim organisation.

The victim profile accepts a business or service name, a daily, monthly, or yearly profit figure, and an estimated hourly profit value associated with the online service being available. This hourly value is intended to capture revenue or business value generated by an online ordering system, QR-code service, booking platform, application endpoint, website, API, or other internet-facing service. Availability-oriented attacks, including DoS, DDoS, LDoS, and LDDoS, use this hourly value to model simulated lost revenue caused by outage, latency, or service degradation. For example, low-rate denial may leave the service reachable but slow, and is therefore modelled as a partial rather than total loss of online-service value.

Sustainability-oriented attacks, including EDoS, DoW, DDoW, and DoA, use interim metered-cost coefficients to model cloud, serverless, and AI resource exposure. These attacks may generate cost even where the service remains available. The simulator therefore separates metered resource cost from availability loss and reports a running simulated exposure value. Each animated packet, invocation, cost event, or token event displays an indicative exposure value so that users can see the financial effect accumulate during the visual simulation.

The mitigation panel was also extended. Each technical control now has a simulated financial effect expressed as a remaining exposure multiplier. Selecting a mitigation control updates the simulation, reduces the future running exposure rate, and displays the estimated exposure avoided by the control. For example, rate limiting, cloud budget guardrails, serverless concurrency caps, and AI token limits reduce future simulated exposure according to attack-specific interim coefficients. These values are educational defaults and are not empirical claims. They are designed to be replaced with measured values once validated research data becomes available.
